\chapter{Спектральный родитель конденсата стрелочной когерентности}
\label{ch:v7-edge-coherence-spectral-parent-gate}

\section{Уточнение открытого разрыва}

Предыдущий гейт использовал сумму радиального квадрата и сырой нормы
антисимметризованного тензора. Такое действие имело правильные минимумы, но
оставалось составным. Теперь проверяется более сильная конструкция:
получаются ли оба ранговых инварианта из обычных следов одного нечётного
оператора, без отдельного проектора на миноры.

Пусть
\begin{equation}
 V=\mathbb C_{\rm copy}^2,\qquad
 W=\mathbb C_{\rm ch}^3,\qquad
 B\in\operatorname{Hom}(W,V).
 \label{eq:v7-coherence-parent-spaces}
\end{equation}
Рассмотрим градуированную цепь
\begin{equation}
 \mathcal H_0=\mathbb C
 \xrightarrow{\ A_B\ }
 \mathcal H_1=V\otimes W^*
 \xrightarrow{\ C_B\ }
 \mathcal H_2=\Lambda^2V\otimes\Lambda^2W^*,
 \qquad
 \dim_{\mathbb C}\mathcal H=(1,6,3).
 \label{eq:v7-coherence-parent-chain}
\end{equation}
Первое ребро и поляризованное второе ребро определены формулами
\begin{equation}
 A_B(1)=\operatorname{vec}B,
 \qquad
 C_B=\frac12d(\Lambda^2)_B.
 \label{eq:v7-coherence-parent-edges}
\end{equation}
Множитель $1/2$ не является весом сектора: для квадратичной карты он
фиксирован тождеством Эйлера
\begin{equation}
 C_BA_B=\Lambda^2B.
 \label{eq:v7-coherence-parent-polarization}
\end{equation}

\section{Один нечётный оператор}

На чётно-нечётно-чётной цепи положим
\begin{equation}
 \mathcal D_B=
 \begin{pmatrix}
 0&A_B^\dagger&0\\
 A_B&0&C_B^\dagger\\
 0&C_B&0
 \end{pmatrix},
 \qquad
 \Gamma=\operatorname{diag}(1,-I_6,I_3).
 \label{eq:v7-coherence-parent-dirac}
\end{equation}
Тогда $\mathcal D_B^\dagger=\mathcal D_B$ и
$\{\Gamma,\mathcal D_B\}=0$. Обозначим
\begin{equation}
 T=\Tr(BB^\dagger),\qquad d=\det(BB^\dagger).
 \label{eq:v7-coherence-parent-invariants}
\end{equation}
Прямой compound-matrix расчёт даёт
\begin{equation}
 \|A_B\|_F^2=T,\qquad
 \|C_B\|_F^2=\frac12T,\qquad
 \Tr(C_B^\dagger C_B)^2=\frac18(T^2-d).
 \label{eq:v7-coherence-parent-edge-traces}
\end{equation}
Поэтому полные незавешенные следы одного оператора равны
\begin{align}
 \Tr\mathcal D_B^2&=3T,\\
 \Tr\mathcal D_B^4&=\frac94T^2+\frac{15}{4}d.
 \label{eq:v7-coherence-parent-spectral-traces}
\end{align}
Положительный детерминант возник из обычного полного следа. Относительная
норма внешнего квадрата и сырой четырёхиндексный тензор больше не вводятся
как отдельные части действия.

\section{Единый спектральный полином}

Для общего положительного массового масштаба $\mu$ определим
\begin{equation}
 \mathcal S_\mu(B)
 =\frac49\left(
 \Tr\mathcal D_B^4-\mu\Tr\mathcal D_B^2+\mu^2
 \right).
 \label{eq:v7-coherence-parent-action}
\end{equation}
Подстановка следов даёт точное разложение
\begin{equation}
 \boxed{
 \mathcal S_\mu(B)
 =\left(T-\frac{2\mu}{3}\right)^2+\frac53d.}
 \label{eq:v7-coherence-parent-reduction}
\end{equation}
Таким образом, коэффициент $5/3$ между радиальным и ранговым каналами
фиксирован геометрией цепи. Параметр $\mu$ задаёт только общий радиальный
масштаб
\begin{equation}
 T_\star=\frac{2\mu}{3};
 \label{eq:v7-coherence-parent-radius}
\end{equation}
нормированные проекторы и ранговая страта от него не зависят.

При $\mu=9/2$ восстанавливается прежняя конвенция $T_\star=3$, но новый
родитель не воспроизводит ручной коэффициент четыре. Он даёт другой,
полностью спектральный потенциал
\begin{equation}
 \mathcal S_{9/2}(B)=(T-3)^2+\frac53\det(BB^\dagger).
 \label{eq:v7-coherence-parent-normalized-action}
\end{equation}
Это различие не мешает точному выбору ранга: поскольку оба слагаемых
неотрицательны,
\begin{equation}
 \mathcal S_{9/2}(B)=0
 \quad\Longleftrightarrow\quad
 T=3,\qquad\rank B=1.
 \label{eq:v7-coherence-parent-zero-set}
\end{equation}

\section{Запуск и полный гессиан}

В нуле квадратичная часть равна $-6T$, поэтому
\begin{equation}
 \Spec\Hess_0\mathcal S_{9/2}
 =\{-12\}^{\times12}.
 \label{eq:v7-coherence-parent-origin-hessian}
\end{equation}
Для представителя
$B_0=\begin{psmallmatrix}\sqrt3&0&0\\0&0&0\end{psmallmatrix}$
полный вещественный спектр равен
\begin{equation}
 \Spec\Hess_{B_0}\mathcal S_{9/2}
 =\{0\}^{\times7}
 \cup\{10\}^{\times4}
 \cup\{24\}.
 \label{eq:v7-coherence-parent-vacuum-hessian}
\end{equation}
Семь нулей снова совпадают с касательным пространством
$(S^3\times S^5)/U(1)$. Четыре моды роста ранга и радиальная мода
положительны. В Real-срезе сигнатуры равны соответственно
$6_-$ в нуле и $3_0+3_+$ в минимуме.

\section{Real-завершение и точная граница}

Сопряжённая цепь с антилинейным обменом даёт стандартное Real-удвоение.
Обычный физический полуслед удаляет удвоение и сохраняет следовые формулы
\eqref{eq:v7-coherence-parent-spectral-traces}; нового веса не возникает.
Действие ковариантно относительно естественного
$U(V)\times U(W)$-действия, поскольку внешняя степень функториальна.

\begin{theorem}[Спектральный родитель ранга-один когерентности]
Трёхузловой нечётный оператор~\eqref{eq:v7-coherence-parent-dirac},
линейный по одному полю $B$, имеет полный спектральный полином
\eqref{eq:v7-coherence-parent-action}. Его относительный коэффициент
$5/3$ фиксирован носителем, нулевой сектор стационарен и неустойчив, а
ненулевые поперечно устойчивые минимумы имеют ранг один и порождают чистые
копийный и канальный проекторы. Ручная норма миноров предыдущего гейта
устранена.
\end{theorem}

Результат является строгим на уровне градуированного спектрального
носителя, но ещё не является конечной спектральной тройкой физической
алгебры. Не доказано, что узлы размерностей $1,6,3$, оба фундаментальных
ребра и их Real-образы возникают в уже допущенном четырёхвершинном
представлении с точным первым порядком и без лишних скаляров.

\begin{nogocriterion}[Запрет преждевременного физического чтения]
Нельзя объявлять цепь новым физическим сектором или завершённым родителем
материи до явной бимодульной типизации всех трёх узлов, проверки первого
порядка, аномалий и совместимости с сохранённым классом $15$. Общий
масштаб $\mu$ также нельзя выдавать за выведенную абсолютную массу.
\end{nogocriterion}

\begin{protocol}
\begin{enumerate}
 \item один нечётный оператор, линейный по $B$: \textbf{построен};
 \item поляризация второго ребра: \textbf{канонична};
 \item полный след второго момента: \textbf{$3T$};
 \item полный след четвёртого момента:
 \textbf{$\frac94T^2+\frac{15}{4}d$};
 \item отдельная норма миноров: \textbf{не используется};
 \item относительный коэффициент детерминанта: \textbf{$5/3$};
 \item нулевой сектор: \textbf{12 отрицательных мод};
 \item минимум: \textbf{ранг один, семь касательных нулей};
 \item поперечный гессиан: \textbf{$10$ четырежды и $24$ один раз};
 \item Real-полуслед: \textbf{не меняет коэффициенты};
 \item строгая конечная спектральная тройка: \textbf{не построена};
 \item итог: \textbf{положительный градуированный спектральный родитель};
 \item следующий гейт:
 \textbf{бимодульный admission цепи $1\to6\to3$}.
\end{enumerate}
\end{protocol}

Машинный сертификат сохранён в
\path{s2t_v7_edge_coherence_spectral_parent_gate_results.json}.

\begin{thebibliography}{9}
\bibitem{v7-parent-spectral-action}
A.~H.~Chamseddine, A.~Connes,
\emph{The Spectral Action Principle}, arXiv:hep-th/9606001.
\bibitem{v7-parent-quiver-morse}
M.~Harada, G.~Wilkin,
\emph{Morse theory of the moment map for representations of quivers},
Geometriae Dedicata 150 (2011), 307--353; arXiv:0807.4734.
\end{thebibliography}